\trace\ is a research prototype and benchmark system, not a clinical device. It is not designed to diagnose, treat, replace professional care, or provide crisis services. All claims in this paper are limited to observable benchmark behavior unless future clinical studies validate other effects.

VSM uses synthetic or curated Vietnamese-student sessions. Such sessions can target route, stage, peer, and safety contracts, but they cannot capture the full complexity of real help-seeking conversations, dialect variation, mixed-language Vietnamese, evolving risk, or long-term therapeutic outcomes. The benchmark also simplifies legitimate clinical variation: more than one supportive response can be appropriate in a real session, while VSM must choose expected stages and peer policies for measurement.

Automated scoring and LLM judging introduce evaluator bias. Deterministic labels may over-penalize acceptable alternatives, while LLM judges may reward fluent responses that hide process failures. The planned human audit is therefore a validity check, not a substitute for clinical evaluation.
